% Chapter 14 — Exceptions
% Source: ../source/split/chapters/ch14-exceptions_p166-p173.pdf
% Original print pages: 171--178
% Status: 已确认

\chapter{异常}
\label{chap:14}

\section*{概述}

\chapref{13}说明了怎样用可变引用扩展纯简单类型 lambda 演算的
简单操作语义，并考察这项扩展对类型规则和类型安全性证明的影响。本章处理
对原计算模型的另一项扩展：抛出和处理异常。

现实编程中，函数经常需要通知调用者：由于某种原因，它无法完成任务。可能是
某项计算涉及除零或算术溢出，字典中缺少查找键，数组索引越界，文件找不到或
无法打开，系统内存耗尽、用户终止进程一类灾难性事件发生，等等。

\taplsecref{11.10}{sec:11.10} 已经看到，有些异常条件可以通过让函数返回变体
（或 option）来表示。但若这些条件确实属于异常情形，我们也许不愿强迫函数的每个
调用者都处理它们发生的可能。更合意的做法可能是：异常条件把控制直接转移到
程序更高层定义的异常处理器；甚至在异常足够罕见、或调用者无论如何也无法
恢复时，直接终止程序。本章先考察后一情形（\taplsecref{14.1}{sec:14.1}），把异常视为
整个程序的终止；再加入捕获异常并从中恢复的机制（\taplsecref{14.2}{sec:14.2}）；
最后细化两种机制，让程序员指定的额外数据可以从异常抛出点传给处理器
（\taplsecref{14.3}{sec:14.3}）。\footnote{本章研究在简单类型 lambda 演算
（\cref{fig:9-1}）上加入多种异常抛出与处理基本构造的系统
（\cref{fig:14-1} 与\cref{fig:14-2}）。第一种扩展的 OCaml 实现是
\texttt{fullerror}；携带值的异常语言（\cref{fig:14-3}）没有官方实现。}

\begin{figure}[!ht]
\centering
\begin{minipage}{0.96\linewidth}
\small
\taplfigureheader{\(\mathsf{error}\)}
  {\(\lambda_{\to}\)（\cref{fig:9-1}）}
\[
\begin{gathered}
\text{\textbf{新增句法形式}}\qquad
\begin{array}{@{}rcll@{}}
t&::=&\cdots\mid\mathsf{error}&\text{项：运行时错误}
\end{array}
\qquad
\begin{gathered}
\text{\textbf{新增求值规则}}\\[-2pt]
\begin{array}{c}
\inferrule*[right=\textsc{E-AppErr1}]
 { }{\mathsf{error}\ t_2\trans\mathsf{error}}
\\[10pt]
\inferrule*[right=\textsc{E-AppErr2}]
 { }{v_1\ \mathsf{error}\trans\mathsf{error}}
\end{array}
\end{gathered}
\\[12pt]
\text{\textbf{新增类型规则}}\qquad
\inferrule*[right=\textsc{T-Error}]
 { }{\Gamma\vdash\mathsf{error}:T}
\end{gathered}
\]
\end{minipage}
\caption{错误}
\label{fig:14-1}
\end{figure}

\begin{figure}[!ht]
\centering
\begin{minipage}{0.96\linewidth}
\small
\taplfigureheader{\(\mathsf{error},\ \mathsf{try}\)}
  {\(\lambda_{\to}\) 并带错误（\cref{fig:14-1}）}
\textbf{新增句法形式}
\[
t::=\cdots\mid\mathsf{try}\ t\ \mathsf{with}\ t
\qquad\text{捕获错误}
\]
\textbf{新增求值规则}
\[
\begin{array}{c@{\qquad}c}
\inferrule*[right=\textsc{E-TryV}]
 { }{\mathsf{try}\ v_1\ \mathsf{with}\ t_2\trans v_1}
&
\inferrule*[right=\textsc{E-TryError}]
 { }{\mathsf{try}\ \mathsf{error}\ \mathsf{with}\ t_2\trans t_2}
\\[12pt]
\inferrule*[right=\textsc{E-Try}]
 {t_1\trans t_1'}
 {\mathsf{try}\ t_1\ \mathsf{with}\ t_2
 \trans\mathsf{try}\ t_1'\ \mathsf{with}\ t_2}
\end{array}
\]
\textbf{新增类型规则}
\[
\begin{array}{c}
\inferrule*[right=\textsc{T-Try}]
 {\Gamma\vdash t_1:T\\\Gamma\vdash t_2:T}
 {\Gamma\vdash\mathsf{try}\ t_1\ \mathsf{with}\ t_2:T}
\end{array}
\]
\end{minipage}
\caption{错误处理}
\label{fig:14-2}
\end{figure}
\FloatBarrier

\translatornote{\cref{fig:14-1} 中的 \(\mathsf{error}\) 表示一种不携带额外信息的
运行时错误，也是异常机制的最简模型。\cref{fig:14-2} 概括了
\taplsecref{14.2}{sec:14.2} 为它加入的处理机制；\cref{fig:14-3} 概括了
\taplsecref{14.3}{sec:14.3} 用 \(\mathsf{raise}\ t\) 对它所作的推广，使异常可以携带数据。}

\begin{figure}[!ht]
\centering
\begin{minipage}{0.99\linewidth}
\footnotesize
\taplfigureheader{\(\mathsf{exceptions}\)}
  {\(\lambda_{\to}\)（\cref{fig:9-1}）}
\textbf{新增句法形式}
\[
t::=\cdots\mid\mathsf{raise}\ t
\mid\mathsf{try}\ t\ \mathsf{with}\ t
\qquad
\begin{array}{l}
\mathsf{raise}\ t\text{：抛出异常}\\
\mathsf{try}\ t\ \mathsf{with}\ t\text{：处理异常}
\end{array}
\]
\textbf{新增求值规则}
\[
\begin{array}{c@{\qquad}c}
\inferrule*[right=\textsc{E-AppRaise1}]
 { }{(\mathsf{raise}\ v_{11})\,t_2\trans\mathsf{raise}\ v_{11}}
&
\inferrule*[right=\textsc{E-AppRaise2}]
 { }{v_1\,(\mathsf{raise}\ v_{21})\trans\mathsf{raise}\ v_{21}}
\\[8pt]
\inferrule*[right=\textsc{E-Raise}]
 {t_1\trans t_1'}{\mathsf{raise}\ t_1\trans\mathsf{raise}\ t_1'}
&
\inferrule*[right=\textsc{E-RaiseRaise}]
 { }{\mathsf{raise}\,(\mathsf{raise}\ v_{11})
      \trans\mathsf{raise}\ v_{11}}
\\[8pt]
\inferrule*[right=\textsc{E-TryV}]
 { }{\mathsf{try}\ v_1\ \mathsf{with}\ t_2\trans v_1}
&
\inferrule*[right=\textsc{E-TryRaise}]
 { }{\mathsf{try}\ (\mathsf{raise}\ v_{11})\ \mathsf{with}\ t_2
      \trans t_2\,v_{11}}
\\[8pt]
\multicolumn{2}{c}{
\inferrule*[right=\textsc{E-Try}]
 {t_1\trans t_1'}
 {\mathsf{try}\ t_1\ \mathsf{with}\ t_2
  \trans\mathsf{try}\ t_1'\ \mathsf{with}\ t_2}}
\end{array}
\]
\textbf{新增类型规则}
\[
\inferrule*[right=\textsc{T-Exn}]
 {\Gamma\vdash t_1:T_{\mathit{exn}}}
 {\Gamma\vdash\mathsf{raise}\ t_1:T}
\qquad
\inferrule*[right=\textsc{T-Try}]
 {\Gamma\vdash t_1:T\\
  \Gamma\vdash t_2:T_{\mathit{exn}}\to T}
 {\Gamma\vdash\mathsf{try}\ t_1\ \mathsf{with}\ t_2:T}
\]
\end{minipage}
\par\medskip
\caption{携带值的异常}
\label{fig:14-3}
\end{figure}
\FloatBarrier

\section{抛出异常}
\label{sec:14.1}

先为简单类型 lambda 演算加入表示异常的最简单机制：项
\(\mathsf{error}\) 一经求值，就完全终止包含它的项的求值。所需扩展见
\cref{fig:14-1}。

为 \(\mathsf{error}\) 编写规则时，主要设计决定是怎样在操作语义中形式化
“异常终止”。这里采用一个简单办法：让 \(\mathsf{error}\) 本身作为程序
终止的结果。规则 \textsc{E-AppErr1} 与 \textsc{E-AppErr2} 刻画这种
行为。\textsc{E-AppErr1} 说明，在试图把应用左端归约为值的过程中遇到
\(\mathsf{error}\) 时，立即让整个应用以 \(\mathsf{error}\) 为结果。
类似地，\textsc{E-AppErr2} 说明，在把应用参数归约为值的过程中遇到错误
时，放弃这个应用并立即产生 \(\mathsf{error}\)。

注意，\(\mathsf{error}\) 只加入项语法，没有加入值语法。这保证
\textsc{E-AppAbs} 与 \textsc{E-AppErr2} 的左端永不重叠。也就是说，
对项
\[
(\lambda x:\tapltype{Nat}.\,0)\ \mathsf{error}
\]
不会产生“究竟执行应用得到 \(0\)，还是终止”的歧义；只有终止这一种可能。
类似地，规则 \textsc{E-AppErr2} 的左端写作
\(v_1\ \mathsf{error}\)，而不是 \(t_1\ \mathsf{error}\)。这表示只有当应用
左端已经求值为值时，右端的 \(\mathsf{error}\) 才会向外传播。因此，在
\[
(\mathsf{fix}\,(\lambda x:\tapltype{Nat}.\,x))\ \mathsf{error}
\]
中，求值器必须先处理应用左端。左端会不断回到自身，始终不能成为值，所以
整个项将一直求值下去，不会传播右端的 \(\mathsf{error}\)。

这项限制与前面“不把 \(\mathsf{error}\) 视为值”的规定共同消除了求值规则
之间的重叠：对于任意项，每一步至多只有一条规则能够决定其下一步结果。因此，
加入 \(\mathsf{error}\) 后，求值关系仍然是确定的。

类型规则 \textsc{T-Error} 也很有意思。我们可能要在任何上下文抛出异常，
所以允许项 \(\mathsf{error}\) 具有任意类型。在
\[
(\lambda x:\tapltype{Bool}.\,x)\ \mathsf{error}
\]
中，它具有类型 \(\tapltype{Bool}\)；在
\[
(\lambda x:\tapltype{Bool}.\,x)\,
(\mathsf{error}\ \tapltrue)
\]
中，它具有类型 \(\tapltype{Bool}\to\tapltype{Bool}\)。

\translatornote{本章只在类型中记录计算正常完成时的结果。一个正常结果为 \(T\)
的计算，即使可能发生 \(\mathsf{error}\)，类型仍然写作 \(T\)，不会改成 Rust 中的
\texttt{Result<T, E>}、\texttt{Option<T>} 或和类型 \(T+\mathit{Error}\)。在这种
设计下，错误向外传播并取代一个类型为 \(T\) 的表达式后，传播得到的
\(\mathsf{error}\) 也必须能够具有类型 \(T\)，这样良类型项经过求值后才能继续保持
良类型。规则 \textsc{T-Error} 因而允许 \(\mathsf{error}\) 具有任意类型。发生错误时，
计算会异常终止，不会产生正常结果；这一点与 Rust 中表示永不返回的 \texttt{!} 类型相近。}

\(\mathsf{error}\) 的类型如此灵活，会给类型检查算法的实现带来困难，因为
语言中每个可赋型项具有唯一类型这一性质（\cref{thm:9.3.3}）不再成立。
可以用多种办法处理。在带子类型化的语言中，可以为 \(\mathsf{error}\)
指派最小类型 \(\tapltype{Bot}\)（见
\taplsecref{15.4}{sec:15.4}），再按需提升为
其他类型。在带参数多态的语言中（见\chapref{23}），可以为它赋予
多态类型 \(\forall X.\,X\)，该类型可以实例化为任何其他类型。两种技巧都
能用一个类型紧凑表示 \(\mathsf{error}\) 无穷多个可能类型。

\begin{exercise}[\(\star\)]
\label{ex:14.1.1}
要求程序员在每处使用 \(\mathsf{error}\) 时标注预期类型，不是更简单吗？
\end{exercise}
\taplauthorsolutionlink{sol:author:14.1.1}{14.1.1}

带异常语言的类型保持性质仍和以往一样：若项具有类型 \(T\)，则它经过一次
单步求值后仍具有类型 \(T\)。进展性质则要稍加细化。它原本说明良类型程序必须求值
为值（或发散）；现在新增了非值范式 \(\mathsf{error}\)，它完全可能成为
良类型程序的求值结果，所以必须在进展性陈述中允许它。

\begin{theorem}[进展性]
\label{thm:14.1.2}
设 \(t\) 是闭的良类型范式。那么，\(t\) 要么是值，要么
\(t=\mathsf{error}\)。
\end{theorem}

\section{处理异常}
\label{sec:14.2}

\(\mathsf{error}\) 的求值规则可以理解为
\termfirst{栈展开}{stack unwinding}：不断丢弃尚未完成的
函数调用，直到错误一路传播到最外层。真实异常语言的实现正是这样工作：
调用栈由一组\termfirst{活动记录}{activation record}组成，每个活动函数调用对应一条；抛出异常会从调用栈中
依次弹出活动记录，直到栈空。

多数异常语言还允许在调用栈中安装异常处理器。抛出异常时，活动记录依次弹出，
直到遇到异常处理器；随后从该处理器继续求值。换言之，异常充当一次非局部
控制转移，其目标是最近安装的异常处理器，也就是调用栈中离它最近的处理器。

\cref{fig:14-2} 汇总的异常处理器形式与 ML、Java 都很相似。表达式
\(\mathsf{try}\ t_1\ \mathsf{with}\ t_2\) 意为：“返回 \(t_1\) 的求值
结果；除非它异常终止，此时改为求值处理器 \(t_2\)。”
\textsc{E-TryV} 说明，\(t_1\) 归约为值 \(v_1\) 后，可以丢掉
\(\mathsf{try}\)，因为已经知道用不到它。\textsc{E-TryError} 则说明，
若 \(t_1\) 的求值结果是 \(\mathsf{error}\)，就用 \(t_2\) 替换整个 try，
从那里继续求值。\textsc{E-Try} 说明，在 \(t_1\) 归约为值或
\(\mathsf{error}\) 之前，应当只继续处理 \(t_1\)，而不动 \(t_2\)。

try 的类型规则直接来自操作语义。整个 try 的结果可能是主表达式 \(t_1\)
的结果，也可能是处理器 \(t_2\) 的结果；只需要求二者具有同一类型 \(T\)，
这也是整个 try 的类型。

类型安全性质及其证明与上一节相比，本质上没有变化。

\section{携带值的异常}
\label{sec:14.3}

\taplsecref{14.1}{sec:14.1} 与 \taplsecref{14.2}{sec:14.2} 的机制让函数可以通知调用者“发生了
某种不寻常的事情”。通常还需要传回额外信息，说明具体发生了哪种事情，因为
处理器应该采取的行动——恢复后重试，或向用户呈现易懂的错误消息——可能
取决于这些信息。

\cref{fig:14-3} 说明怎样扩充基本异常处理构造，让每个异常携带一个值。
这个值的类型写作 \(T_{\mathit{exn}}\)。暂且不确定它的具体性质，稍后讨论
几种选择。

原子项 \(\mathsf{error}\) 被项构造子 \(\mathsf{raise}\ t\) 取代，
\(t\) 是要传给异常处理器的额外信息。try 的语法不变，但
\(\mathsf{try}\ t_1\ \mathsf{with}\ t_2\) 中的处理器 \(t_2\) 现在被
解释为以额外信息为参数的函数。

求值规则 \textsc{E-TryRaise} 实现这种行为：取得 \(t_1\) 中
\(\mathsf{raise}\) 携带的额外信息，再传给处理器 \(t_2\)。
\textsc{E-AppRaise1} 与 \textsc{E-AppRaise2} 像
\cref{fig:14-1} 的 \textsc{E-AppErr1} 与 \textsc{E-AppErr2} 一样，
使异常穿过应用向外传播。不过，这些规则只允许传播额外信息已经是值的异常；
若试图求值的 \(\mathsf{raise}\) 所携信息本身还需要求值，规则会阻塞，
迫使我们先用 \textsc{E-Raise} 求值额外信息。\textsc{E-RaiseRaise}
传播在求值某个异常所要携带的信息时发生的异常。\textsc{E-TryV} 与
\taplsecref{14.2}{sec:14.2} 一样，说明主表达式归约为值后可以丢弃 try。
\textsc{E-Try} 则让求值器继续处理 try 体，直到它变为值或
\(\mathsf{raise}\)。

类型规则反映这些行为变化。在 \textsc{T-Exn} 中，要求额外信息具有类型
\(T_{\mathit{exn}}\)；整个 \(\mathsf{raise}\) 则可以具有上下文所需的任意
类型 \(T\)。在 \textsc{T-Try} 中，检查处理器 \(t_2\) 是一个函数：给它
\(T_{\mathit{exn}}\) 型的额外信息，它产生与 \(t_1\) 相同类型的结果。

最后，考虑 \(T_{\mathit{exn}}\) 的几种选择。

\begin{enumerate}
  \item 可以令 \(T_{\mathit{exn}}=\tapltype{Nat}\)。这对应 Unix 操作
    系统函数等使用的 \texttt{errno} 约定：每个系统调用都返回数字“错误码”，
    其中 \(0\) 表示成功，其他值报告不同的异常条件。

  \item 可以令 \(T_{\mathit{exn}}=\tapltype{String}\)。这样不必查错误码
    表，异常抛出点也可以按需构造更具描述性的消息。额外灵活性的代价是，
    处理器可能需要解析这些字符串，才能知道发生了什么。

  \item 若把 \(T_{\mathit{exn}}\) 定义成变体类型，就能既传递更丰富的异常，
    又避免解析字符串：
    \[
    \begin{aligned}
    T_{\mathit{exn}}=\langle&
      \mathit{divideByZero}:\tapltype{Unit},
      \mathit{overflow}:\tapltype{Unit},\\
      &\mathit{fileNotFound}:\tapltype{String},
      \mathit{fileNotReadable}:\tapltype{String},\ldots\rangle
    \end{aligned}
    \]
    处理器可以用简单的 case 表达式区分异常种类；不同异常也能携带不同类型的
    附加信息。除零等异常无需额外负担；找不到文件则可携带字符串，指出打开
    哪个文件时发生错误。

    这种方案的问题是很不灵活：要求预先固定任意程序可能抛出的全部异常，
    即变体类型 \(T_{\mathit{exn}}\) 的标签集合，程序员无法声明应用特有的
    异常。

  \item 可以进一步把 \(T_{\mathit{exn}}\) 取作
    \termfirst{可扩展变体类型}{extensible variant type}，为用户定义异常
    留出空间。ML 采用这种思路，提供名为 \texttt{exn} 的单一可扩展变体
    类型。\footnote{也可以更进一步，把可扩展变体类型作为一般语言特性；
    ML 设计者选择只把 \texttt{exn} 当作特例。}在当前体系中，ML 声明
    \(\mathsf{exception}\ l\ \mathsf{of}\ T\) 可理解成如下操作：
    保证 \(l\) 不同于 \(T_{\mathit{exn}}\) 中已有的全部标签；从现在开始，
    若声明前的可能变体为 \(l_1:T_1,\ldots,l_n:T_n\)，就令%
    \footnote{异常形式是绑定子，因此需要时总能对它作 \(\alpha\) 变换，
    保证 \(l\) 不同于 \(T_{\mathit{exn}}\) 中已经使用的标签。}
    \[
      T_{\mathit{exn}}=
      \langle l_1:T_1,\ldots,l_n:T_n,l:T\rangle
    \]

    ML 抛出异常的语法是 \(\mathsf{raise}\ l(t)\)，其中 \(l\) 是当前作用域
    定义的异常标签。它可以理解为加标签运算符与简单
    \(\mathsf{raise}\) 的组合：
    \[
    \mathsf{raise}\ l(t)
    \quad\triangleq\quad
    \mathsf{raise}\,
      (\langle l=t\rangle\ \mathsf{as}\ T_{\mathit{exn}})
    \]
    类似地，ML 的 try 可以用简单 try 加 case 去糖：
    \[
    \begin{aligned}
    &\mathsf{try}\ t\ \mathsf{with}\ l(x)\Rightarrow h\\
    &\quad\triangleq\quad
    \mathsf{try}\ t\ \mathsf{with}\
    (\lambda e:T_{\mathit{exn}}.\
      \mathsf{case}\ e\ \mathsf{of}\
      \langle l=x\rangle\Rightarrow h
      \mid\ \_\Rightarrow\mathsf{raise}\ e)
    \end{aligned}
    \]
    case 检查抛出的异常是否带标签 \(l\)。若是，就把异常携带的值绑定到
    \(x\)，求值处理器 \(h\)；若不是，就落入通配分支，重新抛出异常。异常
    会继续传播——途中也许被捕获后又抛出——直到到达愿意处理它的处理器，
    或一路到达最外层并终止整个程序。

  \item Java 使用类而不是可扩展变体来支持用户定义异常。语言提供内建类
    \texttt{Throwable}；它或其任意子类的实例都可以用于
    \texttt{throw}（相当于这里的 \(\mathsf{raise}\)）或
    \texttt{try...catch}（相当于这里的 \(\mathsf{try...with}\)）语句。
    只需定义 \texttt{Throwable} 的新子类，就能声明新异常。

    这种异常机制与 ML 的机制其实密切对应。粗略地说，Java 异常对象的运行时
    表示由一个表示其类的标签——直接对应 ML 的可扩展变体标签——和一条实例
    变量记录——对应标签携带的附加信息——组成。

    Java 异常在几个方面比 ML 更进一步。一是异常标签上存在由子类次序生成的
    自然偏序；异常 \(l\) 的处理器实际会捕获携带 \(l\) 类或其任意子类对象
    的所有异常。二是 Java 区分\tapltranslateditalic{异常}与
    \tapltranslateditalic{错误}：前者是内建类
    \texttt{Exception} 的子类，而 \texttt{Exception} 又是
    \texttt{Throwable} 的子类，应用程序可能希望捕获并尝试从中恢复；后者
    是同样继承自 \texttt{Throwable} 的 \texttt{Error} 子类，表示通常应
    直接终止执行的严重条件。两者的关键区别在类型检查规则：方法必须显式声明
    可能抛出哪些异常，但不必声明可能抛出哪些错误。
\end{enumerate}

\begin{exercise}[\(\star\star\star\)]
\label{ex:14.3.1}
上面对第 4 种方案的说明还只是一个非正式草图。请补全其语法、类型规则和
求值规则，给出一套完整、严格的形式定义。
\end{exercise}
\taplauthorsolutionlink{sol:author:14.3.1}{14.3.1}

\begin{exercise}[\(\star\star\star\star\)]
\label{ex:14.3.2}
上文指出，Java 异常——即 \texttt{Exception} 的子类——比 ML 中的异常以及
\cref{fig:14-3} 所定义的异常系统中的异常受到更严格控制：方法可能抛出的
\tapltranslateditalic{每个异常}都必须在方法
类型中声明。扩展\cref{ex:14.3.1} 的解答，让函数类型除了表示参数与
结果类型，也表示它可能抛出的异常集合。证明你的系统类型安全。
\end{exercise}
\taplauthorsolutionlink{sol:author:14.3.2}{14.3.2}

\begin{exercise}[\(\star\star\star\)]
\label{ex:14.3.3}
许多其他控制构造也能用类似本章的技术形式化。熟悉 Scheme 的
\texttt{call-with-current-continuation}（通常缩写为 \texttt{call/cc}）运算符的读者
（参见 Clinger、Friedman 与 Wand，1985；Kelsey、Clinger 与 Rees，
1998；Dybvig，1996；Friedman、Wand 与 Haynes，2001），可以尝试以类型
\(\tapltype{Cont}\ T\) 为基础构造类型规则；该类型表示期待 \(T\) 型参数的
\(T\)-续延。
\end{exercise}
\taplauthorsolutionlink{sol:author:14.3.3}{14.3.3}
